% ============================================================================
% Preamble for the ~300-Page Structuralist AMS Book Monograph (`paper_300p/`)
% Title: Arakelov-Filtered Automorphic Period Towers and the Generalized
%        Riemann Hypothesis
% Dual-engine compatible: pdfLaTeX (\ifPDFTeX) and XeLaTeX / LuaLaTeX
% ============================================================================
\documentclass[11pt,oneside]{amsbook}

\usepackage{iftex}
\ifPDFTeX
  \usepackage[utf8]{inputenc}
  \usepackage[T1]{fontenc}
  \usepackage{lmodern}
  % Safe Unicode compatibility macros for pdfLaTeX
  \DeclareUnicodeCharacter{00A0}{~}
  \DeclareUnicodeCharacter{00B1}{\ensuremath{\pm}}
  \DeclareUnicodeCharacter{00B7}{\ensuremath{\cdot}}
  \DeclareUnicodeCharacter{00D7}{\ensuremath{\times}}
  \DeclareUnicodeCharacter{2013}{--}
  \DeclareUnicodeCharacter{2014}{---}
  \DeclareUnicodeCharacter{2018}{`}
  \DeclareUnicodeCharacter{2019}{'}
  \DeclareUnicodeCharacter{201C}{``}
  \DeclareUnicodeCharacter{201D}{''}
  \DeclareUnicodeCharacter{2026}{\ldots}
  \DeclareUnicodeCharacter{2032}{\ensuremath{'}}
  \DeclareUnicodeCharacter{2080}{\textsubscript{0}}
  \DeclareUnicodeCharacter{2081}{\textsubscript{1}}
  \DeclareUnicodeCharacter{2082}{\textsubscript{2}}
  \DeclareUnicodeCharacter{2083}{\textsubscript{3}}
  \DeclareUnicodeCharacter{2084}{\textsubscript{4}}
  \DeclareUnicodeCharacter{2085}{\textsubscript{5}}
  \DeclareUnicodeCharacter{2086}{\textsubscript{6}}
  \DeclareUnicodeCharacter{2087}{\textsubscript{7}}
  \DeclareUnicodeCharacter{2088}{\textsubscript{8}}
  \DeclareUnicodeCharacter{2089}{\textsubscript{9}}
  \DeclareUnicodeCharacter{2192}{\ensuremath{\to}}
  \DeclareUnicodeCharacter{21A6}{\ensuremath{\mapsto}}
  \DeclareUnicodeCharacter{2208}{\ensuremath{\in}}
  \DeclareUnicodeCharacter{2209}{\ensuremath{\notin}}
  \DeclareUnicodeCharacter{2212}{\ensuremath{-}}
  \DeclareUnicodeCharacter{221E}{\ensuremath{\infty}}
  \DeclareUnicodeCharacter{2229}{\ensuremath{\cap}}
  \DeclareUnicodeCharacter{222A}{\ensuremath{\cup}}
  \DeclareUnicodeCharacter{2248}{\ensuremath{\approx}}
  \DeclareUnicodeCharacter{224D}{\ensuremath{\asymp}}
  \DeclareUnicodeCharacter{2260}{\ensuremath{\neq}}
  \DeclareUnicodeCharacter{2261}{\ensuremath{\equiv}}
  \DeclareUnicodeCharacter{2264}{\ensuremath{\le}}
  \DeclareUnicodeCharacter{2265}{\ensuremath{\ge}}
  \DeclareUnicodeCharacter{2282}{\ensuremath{\subset}}
  \DeclareUnicodeCharacter{2286}{\ensuremath{\subseteq}}
  \DeclareUnicodeCharacter{2295}{\ensuremath{\oplus}}
  \DeclareUnicodeCharacter{2297}{\ensuremath{\otimes}}
  \DeclareUnicodeCharacter{22C6}{\ensuremath{\star}}
\else
  \usepackage{fontspec}
\fi

\usepackage{microtype}
\usepackage[margin=1in]{geometry}
\usepackage{amsmath,amssymb,amsthm,mathtools,mathrsfs,bm}
\usepackage{enumitem}
\usepackage{booktabs,longtable,array}
\usepackage{xcolor}
\usepackage[colorlinks=true,linkcolor=blue!60!black,citecolor=teal!70!black,urlcolor=blue!70!black]{hyperref}
\usepackage[nameinlink,capitalize]{cleveref}

\allowdisplaybreaks
\setcounter{tocdepth}{2}

% Zero-dependency averaged integral (\fint)
\def\Xint#1{\mathchoice
  {\XXint\displaystyle\textstyle{#1}}%
  {\XXint\textstyle\scriptstyle{#1}}%
  {\XXint\scriptstyle\scriptscriptstyle{#1}}%
  {\XXint\scriptscriptstyle\scriptscriptstyle{#1}}%
  \!\int}
\def\XXint#1#2#3{{\setbox0=\hbox{$#1{#2#3}{\int}$}\vcenter{\hbox{$#2#3$}}\kern-.5\wd0}}
\providecommand{\fint}{\Xint-}

% Theorem environments (numbered within chapters)
\numberwithin{equation}{chapter}
\numberwithin{section}{chapter}

\theoremstyle{plain}
\newtheorem{theorem}{Theorem}[chapter]
\newtheorem{proposition}[theorem]{Proposition}
\newtheorem{lemma}[theorem]{Lemma}
\newtheorem{corollary}[theorem]{Corollary}

\theoremstyle{definition}
\newtheorem{definition}[theorem]{Definition}
\newtheorem{notation}[theorem]{Notation}
\newtheorem{hypothesis}[theorem]{Hypothesis}

\theoremstyle{remark}
\newtheorem{remark}[theorem]{Remark}
\newtheorem{example}[theorem]{Example}

% ----------------------------------------------------------------------------
% Provenance, Literature, and Simplification Macros (verified by audit_300p.py)
% ----------------------------------------------------------------------------
% \leanref{src/Relative/Path.lean:line}{Full.DeclName}
% Always escape underscores as \_ in both arguments.
\newcommand{\leanref}[2]{%
  \par\smallskip\noindent{\footnotesize\textcolor{gray!55!black}{%
  \textbf{[Lean Archive]} \texttt{#1} --- \texttt{#2}}}\par\smallskip%
}

% \litref{Author (Year), Exact Theorem / Section}
% Marks classical published theorems with verified hypotheses.
\newcommand{\litref}[1]{%
  \par\smallskip\noindent{\footnotesize\textcolor{teal!65!black}{%
  \textbf{[Literature]} #1}}\par\smallskip%
}

% \newsimplification{Short Descriptive Title}
% Explicitly flags every newly invented mathematical simplification deviating
% from the Lean archive route; ALWAYS followed by a 100% self-contained proof
% and logged in paper_300p/OPTIMIZATION_LOG.md.
\newcommand{\newsimplification}[1]{%
  \par\smallskip\noindent{\footnotesize\textcolor{violet!75!black}{%
  \textbf{[New Simplification]} #1}}\par\smallskip%
}

% ----------------------------------------------------------------------------
% Standard Mathematical Macros
% ----------------------------------------------------------------------------
\newcommand{\NN}{\mathbb{N}}
\newcommand{\ZZ}{\mathbb{Z}}
\newcommand{\QQ}{\mathbb{Q}}
\newcommand{\RR}{\mathbb{R}}
\newcommand{\CC}{\mathbb{C}}
\newcommand{\FF}{\mathbb{F}}
\newcommand{\AAf}{\mathbb{A}}
\newcommand{\GG}{\mathbb{G}}
\newcommand{\PPb}{\mathbb{P}}
\newcommand{\OO}{\mathcal{O}}
\newcommand{\PP}{\mathfrak{p}}
\newcommand{\QQq}{\mathfrak{q}}
\newcommand{\MM}{\mathfrak{m}}
\newcommand{\DD}{\mathfrak{D}}
\newcommand{\GL}{\mathrm{GL}}
\newcommand{\SL}{\mathrm{SL}}
\newcommand{\PGL}{\mathrm{PGL}}
\newcommand{\GSp}{\mathrm{GSp}}
\newcommand{\Sp}{\mathrm{Sp}}
\newcommand{\SU}{\mathrm{SU}}
\newcommand{\UU}{\mathrm{U}}
\newcommand{\GU}{\mathrm{GU}}
\newcommand{\SO}{\mathrm{SO}}
\newcommand{\Gal}{\mathrm{Gal}}
\newcommand{\Hom}{\mathrm{Hom}}
\newcommand{\End}{\mathrm{End}}
\newcommand{\Aut}{\mathrm{Aut}}
\newcommand{\Ind}{\mathrm{Ind}}
\newcommand{\Frob}{\mathrm{Frob}}
\newcommand{\Spec}{\operatorname{Spec}}
\newcommand{\Res}{\operatorname{Res}}
\newcommand{\Tr}{\operatorname{Tr}}
\newcommand{\Nm}{\operatorname{Nm}}
\newcommand{\ord}{\operatorname{ord}}
\newcommand{\len}{\operatorname{len}}
\newcommand{\rk}{\operatorname{rk}}
\newcommand{\vol}{\operatorname{vol}}
\newcommand{\covol}{\operatorname{covol}}
\newcommand{\disc}{\operatorname{disc}}
\newcommand{\cond}{\operatorname{cond}}
\newcommand{\diag}{\operatorname{diag}}
\newcommand{\sgn}{\operatorname{sgn}}
\newcommand{\im}{\operatorname{im}}
\newcommand{\coker}{\operatorname{coker}}
\newcommand{\rd}{\mathrm{rd}}
\newcommand{\HT}{\mathrm{HT}}
\newcommand{\WD}{\mathrm{WD}}
\newcommand{\rec}{\mathrm{rec}}
\newcommand{\Sha}{\mathrm{Sha}}
\newcommand{\Sel}{\mathrm{Sel}}
\newcommand{\Cl}{\mathrm{Cl}}
\newcommand{\HS}{\mathrm{HS}}
\newcommand{\op}{\mathrm{op}}
\newcommand{\cris}{\mathrm{cris}}
\newcommand{\dR}{\mathrm{dR}}
\newcommand{\pst}{\mathrm{pst}}
\newcommand{\ur}{\mathrm{ur}}
\newcommand{\std}{\mathrm{std}}
\newcommand{\ram}{\mathrm{ram}}
\newcommand{\inert}{\mathrm{inert}}
\newcommand{\spl}{\mathrm{split}}
\newcommand{\sig}{\mathrm{sig}}
\renewcommand{\Re}{\operatorname{Re}}
\renewcommand{\Im}{\operatorname{Im}}
\newcommand{\ASK}{\mathrm{ASK}}
\newcommand{\MAIN}{\mathrm{MAIN}}
\newcommand{\TT}{\mathcal{T}}
